% Research fragment, 2026-10-08.  No change to the main manuscript/state.
% A finite second seed with explicit replacement kernels and unit masses.
\subsection{A repaired finite second double-sieve seed}
\label{sec:replacement-second-seed}

The purpose here is to replace the failed non-switched prefix extension
in Wu's second weighted inequality, not to import its old correction
table.  The argument uses the previously derived moving-residue
signed-box interface and positive finite-prime-product upper bound
$K=681/100$.  Its numerical certificate is separate from these inputs.
The source inequality is Wu 2004, Lemma 4.2; its 21 domains and all
their multiplicities are retained below.

Put
\[
 s=\frac{11}{5},\quad s_0=\frac{227}{50},\quad
 \phi_1=\frac{353}{100},\quad\phi_2=\frac{29}{10},\quad
 \phi_3=\frac{61}{25},\qquad
 a_i=1/\phi_i\ (1\le i\le3),\quad a_0=1/s_0,\quad b=1/s.
\]
Thus $a_0<a_1<a_2<a_3<b$.  Letters $a_i$ in this subsection are
relative endpoints, not the absolute initial sieve threshold.

\paragraph{An exact merge lemma.}
If a squarefree specified prefix $B$ belongs to
$\mathcal D^{\rm well}(D)$, and the distinct residual primes
$q_1>\cdots>q_m$ belong to $\mathcal D^{\rm well}(D/B)$ and do
not divide $B$, their sorted union belongs to
$\mathcal D^{\rm well}(D)$.
Indeed, at a residual prime its square product is at most the
residual square product times $B$.  At a fixed prefix prime $p$,
write $P$ for the product of larger fixed primes and $Q$ for the
product of larger residual primes.  If $Q=1$, use the old prefix
test $Pp^2\le D$.  Otherwise the last factor $q_*$ of $Q$ satisfies
$Qq_*\le D/B$ and $q_*>p$, hence
$PQp^2\le(Pp)(D/B)\le D$.
The lemma does not require the residual primes to lie below the
smallest specified prime.

\paragraph{The finite family on which the seed is proved.}
Retain a specified anchor prime of exponent
$t_0\in[29/200,3/20]$.  Fix a total level
$Q=N^\nu$, strictly below $(1+t_0)/2$ on its anchor box.
Every state used here has a nonnegative marked prime-box prefix
$B=N^{T+o(1)}$, containing that anchor, and satisfies
\begin{equation}\label{eq:second-seed-family}
 B\in\mathcal D^{\rm well}(Q),\quad
 T\ge t_0>2\nu-1,\quad T<\nu,\quad
 \text{every specified prime has exponent at most }\mu=25/128.
\end{equation}
Here the coefficient $\sigma$ is a finite nonnegative combination of
characteristic functions of marked prime boxes, with fixed constant
coefficients; no arbitrary prefix-dependent weight is asserted.
All boundary conditions carry fixed margins before passage to limits.
We only use finite collections of boxes, so their prime exponents
and $L=\nu-T$ have fixed positive lower bounds.  Write
\[
 R=Q/B=N^{L+o(1)},\qquad
 \alpha=\frac{1-T}{\nu-T}\ge2,\qquad
 L\le\frac{171}{400},\qquad
 \frac{L}{s}\le\frac{171}{880}<\mu.
\]
The last strict gap is $7/7040$.
Ordinary residual upper and lower weights are constructed at their
actual residual levels.  The merge lemma puts their marked rough
atoms into full well support.  The largest prime of every admitted
atom is at least the anchor and at most $\mu$, so its single-largest
cap $(1+x)/2$ is at least $(1+t_0)/2>\nu$.
Retain all marked slots, original residual factorials and signed
scalars; normalize each final convolution group by its group-size
factorial, restoring that factorial in the scalar.  Repeated-prime
collisions have the existing rough-square power saving.  Omitting
the finitely many specified primes below a residual cutoff changes
the Mertens factor by $1+O(N^{-\epsilon})$ on each fixed box.
Thus the existing signed atom proof, including its coprime mean and
local-scale moving-residue checks, applies after every indicated merge.

The positive product estimate is also needed for each fixed prime
cutoff $N^\epsilon$, not uniformly as $\epsilon$ tends to zero.
Its extension from the old displayed cutoff is obtained by the same
proof: allow at most $\lfloor1/\epsilon\rfloor$ original prime groups
and choose $12\zeta<\epsilon/2$.  In the double Heath--Brown
decomposition a group without a middle variable must contain a long
variable, because its six short variables could not produce a prime
at least $N^\epsilon$.  Three such groups would give three long
variables with product exceeding the total scale.  One and two groups
use the existing long-variable estimates.  A middle variable uses
the same Type II theorem; the complementary divisor order is still
fixed.  All other losses are powers of $N^{-c\epsilon}$ or fixed
logarithmic powers.  The modulus weights and their certified cap $K$
do not change.  This permits finite positive labelled monomial domains,
not arbitrary arithmetic reweightings.  For the single initial seed,
with $\nu$ sufficiently close to its stated cap,
the cutoff already satisfies $L/s_0>100/1149$; its finite descendants may need this fixed
$\epsilon$ extension.

\paragraph{The 21-term weight and its count check.}
Normalize by
\[
 \Theta_Q=4\operatorname{li}(N)
   \sum_B\frac{\sigma(B)\mathcal C(BN)}
                 {\varphi(B)\log(Q/B)},\qquad
 A(w)=\frac{wF(w)}{2e^\gamma},\quad
 \mathfrak a(w)=\frac{wf(w)}{2e^\gamma}.
\]
Write $\Phi_s$ for the prefix average sifted at $R^{1/s}$.
Wu's inequality is
\[
 5\Phi_s\le\Gamma_1-\Gamma_2-\Gamma_3-\Gamma_4
                       +\sum_{i=5}^{21}\Gamma_i+o(\Theta_Q).
\]
It also has a direct squarefree quotient verification.  On the
$R^{a_0}$-rough support let $A,B,C,D$ count prime factors in the
four successive intervals $[a_0,a_1),[a_1,a_2),[a_2,a_3),[a_3,b)$,
with exponents relative to $R$.  Then
\[
 \Gamma_1=4+\mathbf1_{A=0},\quad
 (\Gamma_2,\Gamma_3,\Gamma_4)=(A+B+C+D,A+B,A+B+C).
\]
The positive terms, pointwise and with $x_+=\max(x,0)$, are
\[
\begin{array}{c|l@{\qquad}c|l}
5&\binom{A+B}{2}&6&AC\\
7&(A-1)_+&8&\mathbf1_{A\ge1}B\\
9&\mathbf1_{A=0}(B+C-2)_+&10&\mathbf1_{A=0,B\ge2}(C+D)\\
11&\mathbf1_{A=0,B=1}(C-1)_+&12&\mathbf1_{A\ge2}D\\
13&\mathbf1_{A=1,B\ge1}(C+D)&14&\mathbf1_{A=1,B=0}(C+D-1)_+\\
15&\mathbf1_{A=0,B=1,C\ge1}D&16&\mathbf1_{A=B=0}(C-3)_+\\
17&\mathbf1_{A=B=0,C\ge3}D&18&\mathbf1_{A=B=0,C=2}(D-1)_+\\
19&\mathbf1_{A=0,B=1,C=0}(D-2)_+&20&\mathbf1_{A=B=0,C=1}(D-3)_+\\
21&\mathbf1_{A=B=C=0}(D-5)_+.&&
\end{array}
\]
These formulas follow because a quotient cutoff at the next-to-last
specified prime fixes all preceding specified primes as the smallest
factors.  Their signed total is nonnegative except at $A=B=C=D=0$,
where it is five.  More explicitly, for $A\ge2$ it is
$(A+B-2)(A+B-3)/2+(A-2)C$; for $A=1,B\ge1$ it is
$(B-1)(B-2)/2$; for $A=1,B=0$ it is $(1-C-D)_+$.
For $A=0,B\ge2$ it is $(B-2)(B-3)/2$.
For $A=0,B=1$, it is zero if $C\ge1$ and $(2-D)_+$ otherwise.
For $A=B=0$, it is zero if $C\ge3$, and respectively
$(1-D)_+,(3-D)_+,(5-D)_+$ when $C=2,1,0$.
This proves the pointwise inequality, independently of a table of
correction functions.  Off rough support all terms vanish; square
factors and intersections with the anchor are the usual power-saving
exceptions.

\paragraph{Non-switched kernels and the repair of $\Gamma_5$.}
The baseline upper term is $4A(s_0)+A(\phi_1)$.
The three negative terms have lower main terms
\[
 B(c)=\int_{a_0}^{c}
 \frac{\log(s_0(1-u)-1)}{u(1-u)}\,du,
 \qquad c=b,a_2,a_3.
\]
Their new prefix prime satisfies $u\le b<1/2$, so its singleton
belongs to residual well support; merge it with the old prefix.
Its lower argument $s_0(1-u)$ lies in $[681/275,177/50]\subset(2,4)$.
All needed ordinary lower atoms therefore belong to the family just
checked.  No unproved positive $h$ value is inserted.

For $\Gamma_5,\ldots,\Gamma_8$ the two additional exponents obey
$u<v$.  Their domains and ordinary upper kernels are
\[
\begin{array}{c|c|c}
i&\text{domain}&\text{ordinary kernel against }du\,dv/(uv)\\\hline
5&a_0<u<v<a_2&(1-u-v)^{-1}\\
6&a_0<u<a_1,\ a_2<v<a_3&(1-u-v)^{-1}\\
7&a_0<u<v<a_1&(1-u-v)^{-1}\\
8&a_0<u<a_1,\ a_1<v<a_2&(1-u-v)^{-1}.
\end{array}
\]
All effective upper parameters are at most $s_0-2=127/50<3$.
If an original parameter is below two, lower the sifting cutoff to
half the residual level exponent; this preserves an upper bound and
the displayed kernel because $wF(w)/(2e^\gamma)=1$ up to three.

The residual pair has full well support precisely when
$2v<1$ and $v+2u<1$.  The first inequality always holds.  The second
holds throughout domains $6,7,8$, since their respective suprema are
\[
 a_3+2a_1<1,\qquad3a_1<1,\qquad a_2+2a_1<1.
\]
In domain 5 keep the ordinary kernel on $v+2u<1$.
On its bad corner
\[
 \mathcal B=\{a_0<u<v<a_2:\ v+2u>1\}
\]
use the positive $K$ count.  Its quotient cutoff is $R^{a_0}$,
its unit quotient has a power-saving product range, and its rough
argument is at least $s_0(2-2a_2)>5/3$ for every $\alpha\ge2$.
Thus $\omega\le3/5$ gives the valid replacement kernel
\begin{equation}\label{eq:second-seed-kernel}
 k_5(u,v)=
 \begin{cases}
 (1-u-v)^{-1},&v+2u<1,\\
 3Ks_0/20,&v+2u>1.
 \end{cases}
\end{equation}
The common normalization is $K/4$ before multiplying by $\omega/a_0$.
The bad domain splits exactly at
$u=(1-a_2)/2$, $u=1/3$, and $u=a_2$: its lower $v$ endpoint is
$1-2u$ on the first interval and $u$ on the second.
This is an actual replacement upper bound, not a subtraction of a
missing old $H$ value.

\paragraph{All positive switch domains and their unit masses.}
Use the notation $[\ell,r]^k_\uparrow$ for $k$ ordered exponents
in that interval; juxtaposition denotes independent successive bands.
The unchanged domains $9,\ldots,21$ are
\[
\begin{array}{c|l@{\qquad}c|l}
9&[a_1,a_3]^3_\uparrow&10&[a_1,a_2]^2_\uparrow[a_2,b]\\
11&[a_1,a_2][a_2,a_3]^2_\uparrow&12&[a_0,a_1]^2_\uparrow[a_3,b]\\
13&[a_0,a_1][a_1,a_2][a_2,b]&14&[a_0,a_1][a_2,b]^2_\uparrow\\
15&[a_1,a_2][a_2,a_3][a_3,b]&16&[a_2,a_3]^4_\uparrow\\
17&[a_2,a_3]^3_\uparrow[a_3,b]&18&[a_2,a_3]^2_\uparrow[a_3,b]^2_\uparrow\\
19&[a_1,a_2][a_3,b]^3_\uparrow&20&[a_2,a_3][a_3,b]^4_\uparrow\\
21&[a_3,b]^6_\uparrow.&&
\end{array}
\]
For $r$ extra primes the quotient cutoff is the next-to-last prime,
of exponent $w=u_{r-1}$.  The normalized unsifted mass is the sum
\begin{equation}\label{eq:second-seed-full-mass}
 \int_{\substack{\boldsymbol u\in\mathcal D_i\,,\ \sum u_j\le\alpha-w}}
 \frac{\omega((\alpha-\sum u_j)/w)}{w\prod u_j}\,d\boldsymbol u
 +\int_{\substack{\boldsymbol u\in\mathcal D_i\,,\ \sum u_j=\alpha}}
             \frac{du_1\cdots du_{r-1}}{\prod u_j}.
\end{equation}
In the surface term eliminate a specified coordinate with derivative
one; it counts quotient $1$ and is not part of the ordinary Buchstab
function.  The formula follows by the prime number theorem and the
rough-number estimate on fixed positive boxes.  Boundary strips can
be removed before the limit; an upper bound for the surface term is
also obtained by counting the eliminated prime last.

For $9\le i\le15$, the rough argument is at least $5/3$ already at
$\alpha=2$, hence for every $\alpha\ge2$.
For example the two smallest corner bounds are $47/25$ and
$3136/1765>5/3$; the other five are larger.
Use $\omega\le3/5$.  For $16\le i\le21$ use $\omega\le1$ and
enlarge the nonunit domain if necessary.  There is no unit mass for
$9\le i\le19$, since every sum is at most $4b=20/11<2$.

Put $\ell_{ij}=\log(\phi_i/\phi_j)$,
$\ell_{01}=\log(s_0/\phi_1)$, and
$\ell_{3b}=\log(\phi_3/s)$.
The following explicit masses $M_i$ majorize
\eqref{eq:second-seed-full-mass} for every $\alpha\ge2$:
\begin{align*}
 M_9&=\tfrac35\{(\phi_1+\phi_3)\log(\phi_1/\phi_3)-2(\phi_1-\phi_3)\},\\
 M_{10}&=\tfrac35\{\phi_1-\phi_2(1+\ell_{12})\}\log(\phi_2/s),\\
 M_{11}&=\tfrac35\ell_{12}\{\phi_2(\ell_{23}-1)+\phi_3\},\\
 M_{12}&=\tfrac35\{s_0-\phi_1(1+\ell_{01})\}\ell_{3b},\\
 M_{13}&=\tfrac35\ell_{01}(\phi_1-\phi_2)\log(\phi_2/s),\\
 M_{14}&=\tfrac35\ell_{01}\{\phi_2(\log(\phi_2/s)-1)+s\},\\
 M_{15}&=\tfrac35\ell_{12}(\phi_2-\phi_3)\ell_{3b},\\
 M_{16}&=\phi_2\ell_{23}^4/24,\qquad
 M_{17}=\phi_2\ell_{23}^3\ell_{3b}/6,\\
 M_{18}&=\phi_3\ell_{23}^2\ell_{3b}^2/4,\qquad
 M_{19}=\phi_3\ell_{12}\ell_{3b}^3/6,\\
 M_{20}&=\phi_3\ell_{23}\ell_{3b}^4/24
                         +\phi_2\ell_{3b}^4/24,\\
 M_{21}&=\phi_3\ell_{3b}^6/720+\phi_3\ell_{3b}^5/120.
\end{align*}
The second summand in $M_{20}$ bounds the unit surface by eliminating
$u_1$, using $1/u_1\le\phi_2$, and dropping its interval restriction.
The other four variables are ordered in $[a_3,b]$.
For $M_{21}$ eliminate $u_6$, use $1/u_6\le\phi_3$, and integrate
the other five ordered variables.  Thus these are uniform mass bounds
even when a descendant has $\alpha>2$.

At the initial value $\alpha=2$ there is a sharper useful check:
domain 21 is empty, and domain 20 has no nonunit quotient because
$a_2+5a_3>2$.  Its unit surface has
\[
 \delta=2-a_2-4a_3=28/1769,\qquad
 M_{20}^{\rm unit}(2)\le
 \frac{\delta^4}{(4!)^2a_2a_3^4}
 =\frac{4802}{428712890625}.
\]
Indeed write its last four coordinates as $a_3+x_i$;
$x_i\ge0$, $\sum x_i\le\delta$, ordered, and the first coordinate
is $a_2+\delta-\sum x_i$.  Its other endpoint restrictions are
automatic here.  The ordered simplex volume is $\delta^4/(4!)^2$.
This sharper initial mass is not used when iterating: the uniform
$M_{20},M_{21}$ above are used instead.

\begin{proposition}[A rigorously bounded finite replacement seed]
\label{prop:replacement-second-seed}
On the finite family \eqref{eq:second-seed-family}, with the fixed
positive-cutoff version of the $K$ input just described, one has
\[
 \Phi_{11/5}\le\{1-J_*+o(1)\}\Theta_Q,
 \qquad J_*=9/500.
\]
This statement uses ordinary lower baselines only and invokes no
previous $H,h$ table or unproved child correction.
\end{proposition}

\begin{proof}
Let $U_i$ for $5\le i\le8$ be the integrals of the displayed ordinary
kernels against $du\,dv/(uv)$, and put
\[
 E_5=\iint_{\mathcal B}
 \left\{\frac{3Ks_0}{20}-\frac1{1-u-v}\right\}\frac{du\,dv}{uv}.
\]
It follows from the established termwise bounds and the common
$K/4$ normalization that
\[
 J_2=\frac15\left\{5-4A(s_0)-A(\phi_1)
       +B(b)+B(a_2)+B(a_3)-\sum_{i=5}^8U_i-E_5
                         -\frac K4\sum_{i=9}^{21}M_i\right\}
\]
is a valid saving.  Every quantity here is evaluated at the same
total level $Q$.

The file \texttt{certify\_replacement\_second\_seed.py} uses 80-bit
outward integer intervals, 32 positive atanh terms with explicit
remainders, and 4096 rational Darboux cells for each one-dimensional
integral.  The double ordinary kernels are integrated once explicitly:
\[
 \int_l^r\frac{dv}{uv(1-u-v)}
  =\frac1{u(1-u)}\log\frac{r(1-u-l)}{l(1-u-r)}.
\]
The split bad domain uses its two exact lower endpoints stated above.
The switch bounds are the explicit logarithmic expressions $M_i$;
there is no floating quadrature input.  The certificate gives
\[
 J_2\ge
 \frac{22417233603428441622629}{1208925819614629174706176}
 >\frac9{500}.
\]
The strict surplus absorbs all previously fixed asymptotic and box
errors.  This proves the assertion under the named analytic inputs.
\end{proof}

\paragraph{A finite propagation, with its lower-child domain retained.}
There is a small fully specified use at the original Buchstab cutoff.
Start with an anchor prefix and upper target argument three.  Buchstab
from arguments three to four produces lower children with arguments
$y\in[2,3]$ and remaining exponent at most $3L/4$.
The ordinary upper baseline at four is licensed by the same merge
argument.  On $2\le y\le5/2$, apply a second Buchstab identity from
$y$ to $16/5$ to that child.  Its ordinary lower baseline is licensed
because all residual primes are at most
$(3L/4)/(16/5)<\mu$.  Its upper grandchildren have arguments
$z\in[y-1,11/5]$.
For $z\le11/5$, monotonicity of the positive sifted count gives the
seed bound at argument $z$ from the proved bound at $11/5$;
the normalized $A$ equals one on this entire interval.
It only lowers the comparison cutoff.

Both newly appended singleton primes have square at most the current
residual level, so the merge lemma preserves full prefix well support.
The first lower-child cutoff is at most $3L/8<\mu$, and all second
appended primes have that same upper bound.  Hence no largest prime
leaves $[t_0,\mu]$.  Total prefix exponents increase, so
$\alpha=(1-T)/(\nu-T)$ increases and remains at least two;
the uniform masses, not the initial $\alpha=2$ masses, are used.
Only finitely many boxes and a fixed positive prime cutoff are needed.

Cancellation of the ordinary $A,\mathfrak a$ Buchstab main terms gives
the explicit defect bounds
\begin{equation}\label{eq:replacement-second-seed-h3}
 h_{\rm child}(y)\ge J_*\log\frac{11/5}{y-1}
 \quad(2\le y\le5/2),\qquad
 H_{\rm anchor}(3)\ge
 J_*\log\frac54\log\frac{22}{15}
 >\frac7{4625}.
\end{equation}
The last inequality uses only
$\log(5/4)>2/9$ and $\log(22/15)>14/37$.
Here $h_{\rm child}$ and $H_{\rm anchor}$ name the bounds on these
particular finite nodes, not an unrestricted old one-variable table.
At the actual anchor cutoff $p$, the argument tends to a value in
$[17/6,171/58]\subset(2,3)$ as the level approaches its cap;
count monotonicity therefore retains
the saving $7/4625$.

If the $L_2$ threshold exceeds $3/20$, the block
$29/200<t<3/20$ is available and yields
\begin{equation}\label{eq:replacement-second-seed-l2-gain}
 \Delta c\ge2\frac7{4625}\log\frac{513}{493}
 >\frac{56}{465275}>0.00012.
\end{equation}
This is a concrete small consequence.  It falls far short of the
approximately $0.094$ gap in the best sampled ordinary counting route.
The repaired finite seed itself is about $0.0185$ under the rigorous
coarse envelopes (the floating exact-$\omega$ probe is about $0.02081$).
It may be used as an input to a larger, explicitly checked finite tree;
it does not authorize restoring the old nine-by-nine matrix unchanged.
If another correction already applies on one of these boxes, take the
stronger valid bound there rather than adding overlapping defects.
No new global Chen coefficient, imbalance exponent, or priority claim
is asserted by this fragment.
